\scanpage{158}
\mspage{75}
et considérons
\[
(257)\qquad V_{\Sigma',\Sigma''}^{\Sigma_{\unsure{1}}}
\overset{\mathrm{déf}}=V_{\Sigma',\Sigma''}\setminus V_{\Sigma',\Sigma''}\cap\Sigma
\]
Prenons \deleted{$\unsure{\Sigma=\varnothing}$} $\Sigma'=\Sigma''$, on trouve $\Sigma'\setminus\Sigma$ \deleted{Pour} (donc $\Sigma'$ si $\Sigma=\varnothing$); prenant $\Sigma''=X$, $\Sigma=\Theta$, on trouve $V^*_{\Sigma',X}$. On peut aussi considérer, en plus des données (256), un
\[
\Sigma'_0\subset\Sigma'
\]
et regarder
\[
(258)\qquad V_{\Sigma'\setminus\Sigma'_0,\Sigma''\setminus\Sigma'_0}^{\Sigma\setminus\Sigma\cap\Sigma'_0},
\]
\deleted{mais} qu'on se gardera de confondre avec
\[
V_{\Sigma',\Sigma''}^{\Sigma\cup\Sigma_0}
\]
P.ex. prenons $\Sigma=\varnothing$, $\Sigma'_0\subset\Sigma'\subset\Sigma''=X$, \inserted{i.e. on a $\Sigma'_0\subset\Sigma'\subset X$} il ne faut pas confondre
\[
V_{\Sigma',X}^{\Sigma_0}\quad\text{et}\quad V_{\Sigma'\setminus\Sigma_0,X\setminus\Sigma_0}
\]
p.ex. si $\Sigma'_0=\Sigma'$, le premier terme est $V_{\Sigma',X}^{\Sigma'}$ (p.ex. si $\Sigma'=\Theta$, c'est $V^*_{\Theta,X}$ \deleted{$\unsure{\Theta}$} qui pour un diviseur lisse est \deleted{\unsure{isomorphe}} \inserted{équivalent} à un fibré en disques épointés sur $\Theta$) alors que le deuxième terme est le topos vide.
\margin{157}
